\section*{4 An explicit polynomial family with high measure}
We now describe the family $\{f_n\}_{n\ge1}$, whose $n$-th member $f_n$ is a polynomial in variables $\mathbf x=\{x_1,\ldots,x_{3m}\}$ and $\mathbf y=\{y_1,\ldots,y_{3m}\}$, where $m\in[n/2,2n]$ will be fixed later.
\[
f_n:=\sum_{\substack{S\subseteq[3m]\\|S|=m}}\left(\prod_{i\in S}y_i\right)\cdot\operatorname{NW}_r(\mathbf x_S),
\]
where $\operatorname{NW}_r$ is a variant of the Nisan-Wigderson design polynomial (introduced in [52]), the construction of which is described later and $r$ is a parameter fixed in this construction. Note that $\{f_n\}_{n\ge1}$ is in VNP. Given a monomial, in order to find its coefficient in $f_n$, we first check if the monomial is multilinear and of degree $m$ in $\mathbf y$ variables. If it is so and $S$ is the set of the indices of the $m$ many $\mathbf y$ variables in the monomial then simply return the coefficient of the part of the monomial in $\mathbf x$ variables in $\operatorname{NW}_r(\mathbf x_S)$--this can be done as the Nisan-Wigderson polynomial family is in VNP.
Let $M_1$ and $M_2$ be as in Section3 and $M=[3m]\setminus(M_1\cup M_2)$. Let $f_1$ be the polynomial computed by the pruned circuit $\mathtt D$, which is obtained from $f=f_n$ by setting the variables $\mathbf x_{\overline M}$ and $\mathbf y_{\overline M}$ to field constants as in Section3.2. Let us now see how $\PSP_{M,k,\ell}(f_1)$ is related to $\dim\langle\pi_{\mathbf x}(\mathbf x_M^\ell\partial_{\mathbf x}^k\operatorname{NW}_r)\rangle$.
\begin{authorclaim}{Lemma14}
Let $f_1$ be as defined above. Then,
$\PSP_{M,k,\ell}(f_1)=\dim\langle\pi_{\mathbf x}(\mathbf x_M^\ell\partial_{\mathbf x}^k\operatorname{NW}_r(\mathbf x_M))\rangle$.
\end{authorclaim}
\begin{proof}
The proof follows easily from the following two observations:
\begin{enumerate}
\item The two operations in $\mathbf y$ variables and the three operations in $\mathbf x$ variables (in the definition of $\PSP_{M,k,\ell}$) commute. That is, we have
\[
\sigma_{\mathbf y_M=1}([\pi_{\mathbf x}(\mathbf x_M^\ell\partial_{\mathbf x}^kf_1)]_m)
=\pi_{\mathbf x}(\mathbf x_M^\ell\partial_{\mathbf x}^k(\sigma_{\mathbf y_M=1}([f_1]_m))).
\]
\item $f_1=(\prod_{i\in M}y_i)\cdot\operatorname{NW}_r(\mathbf x_M)+f'$, where $f'\in\mathbb F[\mathbf x_M,\mathbf y_M]$ and $\deg_{\mathbf y}(f')<m$.
\end{enumerate}
From these observations we have that
\begin{align*}
\PSP_{M,k,\ell}(f_1)
&=\dim\left\langle\sigma_{\mathbf y_M=1}([\pi_{\mathbf x}(\mathbf x_M^\ell\partial_{\mathbf x}^kf_1)]_m)\right\rangle\\
&=\dim\left\langle\pi_{\mathbf x}\left(\mathbf x_M^\ell\partial_{\mathbf x}^k
\left(\sigma_{\mathbf y_M=1}\left[\left(\prod_{i\in M}y_i\right)\operatorname{NW}_r(\mathbf x_M)+f'\right]_m\right)\right)\right\rangle\\
&=\dim\left\langle\pi_{\mathbf x}(\mathbf x_M^\ell\partial_{\mathbf x}^k\operatorname{NW}_r(\mathbf x_M))\right\rangle.
\end{align*}
The last equality follows from the fact that
$\sigma_{\mathbf y_M=1}([(\prod_{i\in M}y_i)\cdot\operatorname{NW}_r(\mathbf x_M)+f']_m)=\operatorname{NW}_r(\mathbf x_M)$.
\end{proof}
\paragraph{Construction of $\operatorname{NW}_r$}
Let $d_x=\lfloor\sqrt n/\ln n\rfloor$. Pick an $\alpha$ such that $d_x\lceil d_x^{1+\alpha}\rceil\le n\le2d_x\lceil d_x^{1+\alpha}\rceil$;this forces $\alpha$ to be $\Theta(\ln\ln n/\ln n)$. Let $q$ be a prime number between $\lceil d_x^{1+\alpha}\rceil$ and $2\lceil d_x^{1+\alpha}\rceil$--such a prime exists [26]--and let $m=d_xq$. Thus, $d_x\lceil d_x^{1+\alpha}\rceil\le m\le2d_x\lceil d_x^{1+\alpha}\rceil$ and hence $n/2\le m\le2n$;moreover, it can be easily verified that $d_x\in[\sqrt m/(2\sqrt2\ln m),2\sqrt2\sqrt m/\ln m]$;both being as required in Section3. Also notice that this means $q=\Theta(\sqrt n\ln n)$. Let $\beta=1/\ln m$ and $r=\lfloor((\alpha+\beta)/(2(1+\alpha)))d_x\rfloor-1$, $\mathbf u=(u_{1,1},\ldots,u_{1,q},\ldots,u_{d_x,1},\ldots,u_{d_x,q})$ and define
\[
\operatorname{NW}_r(\mathbf u):=\sum_{\substack{h(z)\in\mathbb F_q[z]\\\deg(h)\le r}}u_{1,h(1)}\cdots u_{d_x,h(d_x)}.
\]
A lower bound on $\dim\langle\pi_{\mathbf x}(\mathbf x_M^\ell\partial_{\mathbf x}^k\operatorname{NW}_r)\rangle$ was proved in [51,85]. Since their analysis continues to hold for our choice of parameters--which only slightly differ from the parameters in [85]--we omit the proof of the following theorem. Moreover, while they prove this lower bound over fields of characteristic zero, the same proof also works if the characteristic is greater than
$q^{(r+1)\cdot\min\{\binom mk\binom m\ell,\binom m{\ell-d_x-k}\}}$.
\begin{authorclaim}{Theorem15 (Lemma5.2 of [51],Lemma4.1 of [85])}
\[
\dim\left\langle\pi_{\mathbf x}(\mathbf x_M^\ell\partial_{\mathbf x}^k\operatorname{NW}_r(\mathbf x_M))\right\rangle
\ge\frac1{m^{O(1)}}\min\left\{\frac1{4^k}\binom m\ell\binom mk,\binom m{\ell+d_x-k}\right\}.
\]
\end{authorclaim}
Hence, from Lemma14 and Theorem15 we get
\begin{authorclaim}{Lemma16}
$\PSP_{M,k,\ell}(f_1)\ge m^{-O(1)}\min\{4^{-k}\binom m\ell\binom mk,\binom m{\ell+d_x-k}\}$.
\end{authorclaim}
